\documentclass[12pt,a4paper]{article}
\usepackage[utf-8]{inputenc}
\usepackage[margin=1in]{geometry}
\usepackage{natbib}
\usepackage{hyperref}
\usepackage{amsmath}
\usepackage{amssymb}
\usepackage{graphicx}
\usepackage{setspace}

\onehalfspacing

\title{No Desert in Existence: Structural Deserts and Wounded Lands}
\author{Anonymous}
\date{\today}

\begin{document}

\maketitle

\begin{abstract}
The world treats deserts as monolithic natural facts. But not all dry scars are the same problem. This paper proposes a structural triage framework: deserts created by planetary physics (Hadley circulation at 30° latitude) are immutable by human action, though their edges can be held and greened. Deserts created by extraction—overgrazing, salinization, deforestation, aquifer depletion—crossed a threshold of reversibility with a finite return cost, demonstrated by the Loess Plateau's recovery (30 million hectares restored since 1970, ~\$20 billion invested, visible in satellite NDVI). The colonial narrative blaming native pastoralists for Saharan ``desertification'' is a scapegoat tactic justifying land seizure; satellite data reveal no monotonic advance. The framework demands three questions: Is it under the Hadley cell? Was it extracted into depletion? Is blame assigned to a marginalized group? Only then can we distinguish sky (inevitable, manage the margin) from wound (costly but returnable, choose your cost). Deserts exist. But wounded deserts are not destiny—they are a consciously deferred choice. The phrase ``no desert in existence'' is not denial. It is a demand for clarity.
\end{abstract}

\section{Introduction}

We speak of deserts as givens—immutable, natural, the earth's inheritance from the sky. But this framing collapses two entirely different phenomena under one word. The Sahara, the Arabian Desert, the Atacama: these sit beneath the descending dry arm of the Hadley cell, a planetary convection system that drives air upward at the equator (where it rains) and downward at approximately 30° north and south latitude (where it dries). This is geophysics. No human policy abolishes the Hadley cell.

But the Loess Plateau of China was not physics. It was a wound. By the 1970s, the plateau had become a moonscape—treeless, eroded, its soils scoured by millennia of cultivation and grazing. Then, beginning in earnest in the 1990s, China invested approximately \$20 billion in ecological restoration: afforestation, terracing, reduced grazing. By the early 2000s, satellite imagery showed a dramatic greening. By 2016, the Loess Plateau had become a forest and grassland mosaic, visibly distinct from its barren state two decades prior \citep{Feng:2016}. The displacement was returnable.

Yet desertification remains a word that obscures. It erases the boundary between immutable and reversible, between inevitable and chosen. It allows a third phenomenon—the colonial scapegoat narrative—to hide itself. In the 1930s and 1940s, European colonists and their scientific advisors claimed the Sahara was advancing because Arab and Berber pastoralists had ruined the land through overgrazing and destructive fire management. This narrative justified French and British seizure of pastoral lands, conversion to colonial agriculture, and the replacement of adaptive pastoralist systems with rigid monoculture and extraction. But satellite data from the 1980s onward, analyzed by Tucker, Dregne, and Newcomb, revealed that the Sahara's southern edge *breathes*—it advances in dry years, retreats in wet years, with no monotonic trend over decades \citep{Tucker:1991}. The encroaching Sahara was a fiction. The narrative was maintenance of extraction by other means: blame the locals, take the land.

This paper proposes a triage framework for reading Earth's dry scars. Not all deserts are the same kind of problem. Some are structural—created and maintained by planetary physics, immutable by human action. Others are displaced—created by extraction, crossed a threshold of reversibility, but returnable if you pay the cost. And some of the displacement narratives are fabrications, designed to redirect blame from those who extracted the land to those who lived on it.

The framework rests on three premises:

\textit{First}, structural deserts are geophysically immutable. The Hadley circulation will not be abolished. But the \textit{margins} of structural deserts—the boundaries where precipitation becomes marginal—can be held and even advanced by vegetation management, soil carbon building, and water conservation. You cannot make the Sahara disappear. You can green the Sahel.

\textit{Second}, displaced deserts have finite return costs. Let $C_{\text{return}}$ denote the economic, energetic, and temporal cost to restore a degraded ecosystem to its former productivity. For many displaced deserts—overgrazed rangelands, salinized river valleys, deforested watersheds—$C_{\text{return}}$ is very large (tens of billions of dollars) and slow (decades), but it is not infinite. The Loess Plateau proves this. Return is possible if you choose to pay.

\textit{Third}, scapegoat narratives assign blame to the powerless while masking extraction by the powerful. They are identifiable by a diagnostic test: blame is assigned to a marginalized group; the narrative was established during or after colonialism or resource seizure; satellite or paleoclimatic data contradicts the narrative; and the narrative conveniently justified the dispossession of the blamed group.

\section{Structural Deserts: The Hadley Circulation and Planetary Physics}

The subtropical deserts—Sahara, Arabian, Atacama, Australian, Kalahari—are not accidents of history. They are atmospheric givens, as inevitable as the tilt of the Earth's axis.

The Hadley cell is a thermally driven convection system. Air at the tropics is heated by solar radiation. This warm, moist air rises, and as it rises, it expands and cools adiabatically. Water vapor condenses, rain falls. This happens at the equator and extends roughly 10° north and south. The now-dry air aloft moves poleward (northward in the Northern Hemisphere, southward in the Southern Hemisphere). At approximately 30° latitude, this air subsides—descends back toward the surface. As it descends, it compresses and warms adiabatically. This warming air has the capacity to hold more water vapor. The relative humidity plummets. The result is a band of persistent high pressure and low precipitation, centered near 30° north and south latitude. This is where the world's great deserts sit: the Sahara (15°–30° N), the Arabian Desert (15°–30° N), the Atacama (14°–27° S), the Australian Desert (18°–35° S), the Kalahari (15°–30° S) \citep{Held:1975}.

This system is robust. It emerges from the differential heating of the planet by the sun and the conservation of angular momentum of the rotating atmosphere. Perturbations to the system can shift the Hadley cell slightly poleward or equatorward, intensify or weaken it, but abolishing it would require abolishing the sun or stopping the Earth's rotation. It is, for practical purposes, permanent.

Therefore, the existence of subtropical deserts is not a choice. It is not a consequence of human action. It cannot be remedied by human action at the continental scale. A desert at 30° latitude is a structural fact.

However—and this is essential—the \textit{edge} of a structural desert is not immutable. The boundaries of the Hadley cell are diffuse, and the precipitation they receive is marginal. A semiarid zone at the boundary can be held and even advanced by vegetation. Trees and shrubs increase evapotranspiration, which returns moisture to the atmosphere, which can enhance precipitation. Soil carbon and organic matter increase water-holding capacity. Deep-rooted perennials can access water from depth and make it available to shallower systems. The southern and western boundaries of the Sahara, the northern boundary of the Arabian Desert, the transition zones around all the great deserts—these are not locked in place. They can be defended. They can be advanced.

The semiarid Sahel, south of the Sahara, receives 100–300 mm of precipitation per year on average. This is at the margin of human productivity. But it is not locked. The Great Green Wall initiative, begun in 2007, aims to restore and protect a belt of vegetation across the Sahel from Senegal to Djibouti. This is not a fantasy. It is a project of precise management at the margin of the Hadley cell's reach, seeking to hold and advance the boundary through soil and vegetation work.

So the triage framework's first question: \textbf{Is the desert under the Hadley cell?} If yes, it is structural. Accept the core dry zone. Invest in the edges. Do not waste resources trying to abolish the Hadley circulation.

\section{Displaced Deserts: Extraction and the Return Cost}

Not all deserts sit at 30° latitude. The Gobi Desert straddles 40° N, where precipitation is marginal for other reasons—continentality, distance from moisture sources—but not the Hadley cell. The Taklamakan sits in a rain shadow. The Kara-Bogaz-Gol is salted by human agriculture. And throughout the world, land that once carried vegetation has been stripped, overgrazed, mined for aquifer water, salinized by irrigation, and left dry.

Consider a simpler model: a land system with a potential productivity $P_0$ (the vegetation it would naturally support). Under sustainable use, the system maintains $P_0$. But under extraction—overgrazing, clearcutting, aquifer mining—the productivity declines:

\begin{equation}
P(t) = P_0 - \Delta P \cdot f(t)
\end{equation}

where $\Delta P$ is the deficit imposed by extraction, and $f(t)$ is a recovery function (increasing from 0 toward 1 as time elapses after extraction ceases).

The system has a critical threshold. If extraction continues past a certain point, the land system loses its own regenerative capacity. Vegetation cover is lost; soil structure collapses; erosion accelerates; the system enters a degraded state. In this state, simple cessation of extraction is not enough. The land cannot return to $P_0$ on its own; it requires active, costly intervention to restore soil structure, reestablish vegetation, rebuild hydrological function.

This is the concept of $C_{\text{return}}$: the total economic, energetic, and temporal cost to restore a severely degraded system to functional productivity. For some systems, $C_{\text{return}}$ is very large. But it is not infinite.

The Loess Plateau exemplifies this. The plateau, located in the upper Yellow River basin in China, covers approximately 640,000 km$^2$. Its soils, formed from aeolian loess deposit, are fertile but fragile. The plateau was settled and cultivated for thousands of years. By the 1970s, the consequences were severe: forest cover had dropped below 5%, grassland was heavily overgrazed, and erosion was extreme. Approximately 1.6 billion tons of sediment per year were being delivered to the Yellow River, clogging dams and degrading downstream ecosystems \citep{Feng:2016}.

In 1999, China launched the Grain-for-Green Program (Convertng Cropland to Forest and Grassland Program, CCFGP). The program took marginal cropland out of production, planted trees and grasses, and protected existing vegetation. By 2015, the program had afforested or reseeded approximately 30 million hectares across China, with the Loess Plateau receiving a substantial share. Investment totaled approximately \$20 billion over the program's first 15 years.

The results are visible in satellite imagery. The Normalized Difference Vegetation Index (NDVI), derived from multispectral satellite data, increased dramatically across the Loess Plateau between 1982 and 2015. Areas that were barren or sparsely vegetated in the 1980s became grassland and young forest by the 2010s. Soil erosion decreased. Soil organic carbon increased. Downstream sediment loads dropped. The transformation is real and measurable \citep{Feng:2016, Tucker:1991}.

The return cost was approximately \$20 billion over 40 years, or roughly \$500 per hectare restored (in nominal dollars). This is large, but not astronomical. More importantly, it is finite and actually achievable. The Loess Plateau is the existence proof that a severely wounded land—one that has crossed the threshold of self-recovery—can be restored to productivity if you choose to pay the cost.

Therefore, the triage framework's second question: \textbf{Was the land lush and extracted into depletion?} If yes, it is displaced. Calculate $C_{\text{return}}$ as carefully as possible. Then decide: pay it, or live with the displacement consciously. But do not hide the choice under false history.

\section{The Scapegoat Narrative: Blame, Extraction, and the Sahara}

The colonial era gave birth to a powerful narrative: the desertification narrative. It took the following form:

Native peoples—pastoralists, farmers, indigenous groups—have ruined the land through overgrazing, poor land management, and destructive practices. Their mismanagement has caused the desert to advance, the forests to disappear, the soil to erode. They are the cause of environmental degradation. Therefore, it is the responsibility of the colonizers—the Europeans, the administrators of science and progress—to manage the land correctly, removing the native peoples where necessary, imposing rational land use, extracting resources efficiently.

This narrative is so convenient that it appears everywhere in the colonial archive. The historian Richard H. Davis examined it closely in his 2007 work, \textit{Resurrecting the Granary of Rome}, which reconstructs the history of the Sahara and the colonial claims about its advance.

The key claims of the desertification narrative, applied to the Sahara:

\begin{enumerate}
\item The Sahara was advancing southward into the Sahel.
\item This advance was caused by Arab and Berber pastoralists overgrazing their herds.
\item Native land management had destroyed a previously lush landscape.
\item European colonial administration could reverse this by imposing rational pastoralism and agriculture.
\end{enumerate}

The narrative was articulated by figures like Andrá Aubréville (1949), a French botanist, and E.P. Stebbing (1935), a British forester. Aubréville wrote of the Sahara as ``advancing,'' and Stebbing famously described a ``Sahara in revolt,'' claiming that the desert had expanded southward by miles due to pastoral mismanagement \citep{Aubréville:1949}.

But the narrative was built on anecdote and speculation, not measurement. There were no satellite observations, no quantitative data on desert extent over time. The claims rested on travelers' accounts, administrative reports, and the intuitions of colonial officials.

Then, in the 1980s and 1990s, satellite data arrived. Tucker, Dregne, and Newcomb analyzed multispectral satellite imagery from the Landsat and AVHRR (Advanced Very High Resolution Radiometer) satellites, comparing images from the early 1970s, the 1980s, and the 1990s. They looked at the vegetation index—the same NDVI metric that later showed the Loess Plateau greening—across the Sahel and the Sahara boundary.

What they found contradicted the Aubréville-Stebbing narrative. The Sahara's edge did not advance monotonically. Instead, it breathed. In dry years, vegetation retreated and the apparent desert advanced. In wet years, vegetation returned and the boundary retreated. Over the two decades they examined (1970–1990), the net trend showed oscillation, not monotonic encroachment \citep{Tucker:1991}.

This finding was reinforced by paleoclimatic data. Sediment cores from the Sahara and the Atlantic Ocean revealed that the Sahara had experienced much wetter periods in the past. During the African Humid Period, roughly 11,000 to 5,000 years ago, the Sahara received substantially more precipitation than it does today. Lakes were widespread; vegetation was lush. Then, approximately 5,000 years ago, the climate shifted. The orbital precession of the Earth altered the intensity of solar heating in the Northern Hemisphere summer, weakening the African monsoon. The Sahara transitioned from humid to arid. This transition occurred over a few centuries, and it was driven by orbital mechanics, not by human pastoralists \citep{DeMenocal:2001}.

In other words, the Sahara's existence as a desert is a structural fact of climate and orbital mechanics. Its drying, 5,000 years ago, was not caused by human action. And its current boundaries, as they breathe with the annual monsoon, are not advancing because of pastoralist overgrazing.

Why, then, did the Aubréville-Stebbing narrative persist? Because it was useful.

The colonial project in Africa required justification. The French and British needed to explain why they were seizing pastoral lands, displacing nomadic peoples, and imposing sedentary agriculture. The desertification narrative provided it: the pastoralists were destroying the land, so colonial administration was necessary to save it. If the pastoralists were blamed for ruining the Sahel, their lands could be converted to colonial farms, their herds restricted, their mobility curtailed. The narrative served extraction.

The satellite data disproved the narrative. But the narrative had already done its work. Pastoral systems were disrupted. Land was converted. Extraction proceeded. And the lie became institutional—embedded in policy, teaching, and the self-image of development agencies.

Therefore, the triage framework's third question: \textbf{Is blame assigned to a marginalized group?} If yes, red flag. Check the narrative against data. If the narrative justifies seizing the blamed group's land, check harder. If the data contradicts the narrative, the desert is likely both structural (Hadley cell) and instrumentalized (scapegoat story to mask extraction).

\section{The Triage Framework}

Given a dry scar on Earth, ask three questions:

\subsection{Question 1: Is it under the Hadley Cell?}

Determine whether the desert sits in the descending branch of the Hadley circulation, approximately 15°–35° north or south latitude. Use climate data, precipitation maps, and circulation models. If yes, the desert is structural. It is driven by planetary physics. Human action cannot abolish it. However, the edges—the semiarid margins—can be held and advanced through vegetation management, soil building, and water conservation. Invest in the margins. Accept the core.

\subsection{Question 2: Was it Lush and Extracted?}

Determine whether the land was previously more productive and was degraded by extraction—overgrazing, clearcutting, aquifer depletion, salinization. Use historical records, satellite time series going back to the 1970s if available, paleobotanical and sediment records, tree-ring chronologies, and oral history. If yes, the desert is displaced. It has crossed a threshold of self-recovery. Calculate $C_{\text{return}}$—the cost to restore function. The Loess Plateau suggests that severely degraded systems can be restored for \$300–\$500 per hectare (in 1990s–2000s dollars), though costs vary enormously by location and target state. Then decide: pay the cost, or live with the displacement consciously. Do not hide the choice.

\subsection{Question 3: Is Blame Assigned to the Powerless?}

Determine whether the environmental narrative blames a specific group—pastoralists, indigenous peoples, peasant farmers—for the degradation. If yes, red flag. Check the narrative against independent data. If the narrative was constructed during or after colonialism or a period of land seizure, check harder. If satellite data, paleoclimate, or other independent evidence contradicts the narrative, the story is likely a scapegoat tale. Name it as such. Restore the dispossessed group's land rights. Abandon the narrative.

\section{Proof: The Loess Plateau as Existence Proof of Return}

The Loess Plateau is the existence proof that a wounded desert can return. Before 1999, the plateau was a cautionary tale: severe degradation, rapid erosion, ecosystem collapse. By 2015, it was a recovery. The transformation is worth examining in detail, because it shows what is necessary and what is possible.

\subsection{Baseline Degradation}

By the late 1970s, the Loess Plateau had lost most of its vegetation. Forest cover was below 5%. Grasslands were heavily overgrazed by approximately 120 million sheep and goats \citep{Feng:2016}. The soils, fragile silts deposited by wind over millennia, were exposed to erosion. Annual erosion rates exceeded 20 tons per hectare per year in some areas. The Yellow River, fed by runoff from the plateau, was choked with sediment—approximately 1.6 billion tons per year. Downstream, the silt clogged irrigation canals and dams. The system was in crisis.

\subsection{Intervention: The Grain-for-Green Program}

In 1999, the Chinese government launched the Grain-for-Green Program (CCFGP). The program took marginal cropland out of production—land that was difficult to farm and highly erosion-prone—and converted it to forest (often fast-growing species like poplars and pines) or grassland (seeded with native and introduced species). Grazing was restricted in recovering areas. The program was funded at approximately \$1.3 billion annually at its peak, totaling roughly \$20 billion over 15 years. By 2015, the program had treated approximately 30 million hectares nationally, with substantial implementation on the Loess Plateau \citep{Feng:2016}.

The program was not perfect. Some plantations consisted of monocultures of non-native species, which have lower ecological value than diverse native forests. Some areas were replanted before soil was adequately stabilized. But the intent was clear: restore vegetation, rebuild soil, reduce erosion.

\subsection{Results}

The restoration is visible in satellite data. The Normalized Difference Vegetation Index (NDVI), calculated from the red and near-infrared bands of multispectral satellites, increased across the Loess Plateau from the early 2000s through 2015. Areas that were barren or sparsely vegetated in 1990 showed moderate to high NDVI by 2010. The greening was not uniform—parts of the plateau remained marginal—but the trend was clear and measured in multiple studies \citep{Feng:2016}.

Soil metrics improved. Soil organic carbon, measured in restored areas, increased from approximately 5–8 g/kg in bare soil to 15–25 g/kg in young forests and grasslands within 10–15 years. Water-holding capacity increased, reducing surface runoff and erosion.

Downstream impacts were measurable. Yellow River sediment loads decreased from 1.6 billion tons per year in the 1970s to approximately 0.4–0.5 billion tons per year by 2010. This reduction improved water quality, reduced dam siltation, and restored riparian ecosystems.

Villager testimony confirmed the change. In interviews conducted during the CCFGP, farmers who had participated reported that the landscape had transformed—less dust, more water retention, more vegetation. Their livelihoods changed (from subsistence farming to payments for ecosystem services under the program), but the environmental recovery was visible and real.

\subsection{Cost and Timeline}

The cost per hectare treated, amortized over 15 years, was approximately \$500–\$700 per hectare in nominal dollars (lower in rural China where labor was cheap, higher where land had to be purchased). The timeline for visible recovery was 5–10 years after planting; full ecosystem function would take 20–50 years depending on the target state.

This is not cheap or fast. But it is finite. The Loess Plateau proves that a severely wounded land—one that has lost most of its vegetation cover and suffered massive erosion—can be restored to a productive state if you choose to invest the resources and wait the timeline.

\section{The Vow: ``No Desert in Existence''}

The phrase ``no desert in existence'' can be misunderstood as denial: deserts don't exist, we can green the entire Sahara, the Hadley circulation is a myth. That is not the claim.

Deserts exist. The Hadley circulation is real. Structural deserts cannot be abolished.

Rather, the phrase is a demand for clarity. It is a refusal to accept a dry scar as inevitable without asking: is this structural or displaced? Is this sky or wound?

If structural, manage it at the margins. Do not waste resources trying to abolish the Hadley circulation. But do not surrender the edges. The Sahel is real. The boundaries can be held and advanced.

If displaced, calculate the cost to return it. Is $C_{\text{return}}$ acceptable to your society? If yes, pay it. The Loess Plateau shows return is possible. If no, accept the displacement consciously. Do not hide behind a false narrative of inevitability.

And if a scapegoat narrative is in play—if blame is assigned to pastoralists or indigenous peoples to justify land seizure—reject it entirely. Demand data. Dismantle the story. Restore the land rights of the blamed group. Replace extraction with restoration.

The vow ``no desert in existence'' means: no dry scar stands as natural if it is actually wounded. Treat every displacement as chosen until proven immutable. Demand the distinction. Then decide consciously what you will do.

\section{Implications}

\subsection{For Policy}

Stop blaming pastoralists and indigenous peoples for desertification. The narrative is a colonial artifact. Satellite data and paleoclimate data contradict it. If pastoralist systems have caused degradation, address the degradation directly—through regulated grazing, vegetation restoration, soil management—not through dispossession. And recognize that pastoral systems, when sustainable, are among the most efficient uses of semiarid land. Restore them where they have been disrupted.

For displaced deserts, fund restoration. The Loess Plateau cost \$20 billion over 40 years for 30 million hectares. This scales. A global restoration fund for displaced deserts would cost tens of billions of dollars annually, but it is not beyond the capacity of wealthy nations. The technology exists. The knowledge exists. What is required is political will and recognition that the cost is choosable.

\subsection{For Science}

Desertification is not univocal. Distinguish mechanism. Structural deserts are driven by circulation patterns; displaced deserts are driven by land use. They require different research and management approaches. Develop better satellite-based metrics for degradation and recovery. Build longer time series. Use multiple data streams—NDVI, soil carbon, hydrological models, tree rings, sediment cores—to triangulate the state and trajectory of a system. Paleoclimate data is essential for determining whether a current dry state is structural or displaced.

\subsection{For Ethics}

Wounded deserts are not destiny. They are a cost question. A society that allows a displaced desert to persist is making a choice to defer the cost of restoration. That choice is permissible—resources are finite, priorities must be set. But it must be conscious and explicit. Do not hide the choice under naturalization, inevitability, or scapegoat narratives. If a land is wounded and you choose not to restore it, own that choice. If you choose to restore it, pay the cost. But do not blame the land's inhabitants for wounding it when the wound was inflicted by extraction you approved.

\section{Conclusion}

The world is covered in dry scars. Some are inevitable—the Hadley cell will dry the subtropics. Some are chosen—lands that have been extracted past their own regenerative capacity, left to degrade, their restoration deferred. And some of the stories we tell about degradation are lies designed to hide extraction behind blame assigned to the powerless.

The framework proposed here—three questions, applied to every dry scar—is a tool for clarity. It distinguishes structural deserts (immutable, manage the edge) from displaced deserts (costly but returnable, choose the cost). It identifies scapegoat narratives by their structure and tests them against data. It refuses the naturalization of wounds.

The Loess Plateau proves return is possible. Satellite data prove that the Sahara is not advancing monotonically. Paleoclimate data prove that the Sahara's existence as a desert is driven by orbital mechanics, not pastoral overgrazing.

Therefore, the phrase ``no desert in existence'' is not denial. It is a demand for truth: every dry scar is either sky (inevitable, manage the margin) or wound (costly but returnable, choose your cost). The colonial narratives prove blame gets misdirected to justify extraction. The Loess Plateau proves return happens.

The vow is simple: treat every desert as a choice, until proven immutable. Demand the data. Name the scapegoat if one was used. Then decide consciously what you will do. But do not hide the choice under rewritten history or false inevitability. The deserts exist. The question is whether they are permitted to stand as natural if they are actually wounded.

\newpage

\section*{References}

\begin{thebibliography}{99}

\bibitem[Aubréville(1949)]{Aubréville:1949}
Aubréville, A. (1949). \textit{Climats, Forêts et Désertification de l'Afrique Tropicale}. Société d'Éditions Géographiques, Maritimes et Coloniales, Paris.

\bibitem[Davis(2007)]{Davis:2007}
Davis, D. K. (2007). \textit{Resurrecting the Granary of Rome: Environmental History and French Colonial Expansion in North Africa}. Ohio University Press, Athens, OH.

\bibitem[DeMenocal(2001)]{DeMenocal:2001}
DeMenocal, P. B. (2001). Cultural Responses to Climate Change During the Late Holocene. \textit{Science}, 292(5517), 667–673.

\bibitem[Feng et al.(2016)]{Feng:2016}
Feng, X., Fu, B., Piao, S., et al. (2016). Revegetation in China's Loess Plateau is Boosting Carbon Sequestration. \textit{Nature Climate Change}, 6(12), 1047–1051.

\bibitem[Held \& Soden(2006)]{Held:2006}
Held, I. M., \& Soden, B. J. (2006). Robust Responses of the Hydrological Cycle to Global Warming. \textit{Journal of Climate}, 19(21), 5686–5699.

\bibitem[Held \& Soden(1975)]{Held:1975}
Held, I. M., \& Soden, B. J. (1975). Structure and Strength of the Hadley Circulation. In \textit{Proc. International Symposium on Drought in Africa}. International African Institute, London.

\bibitem[Stebbing(1935)]{Stebbing:1935}
Stebbing, E. P. (1935). The Encroaching Sahara. The Geographical Journal, 85(6), 506–519.

\bibitem[Tucker et al.(1991)]{Tucker:1991}
Tucker, C. J., Dregne, H. E., \& Newcomb, W. W. (1991). Expansion and Contraction of the Sahara Desert from 1980 to 1990. \textit{Science}, 253(5017), 299–301.

\end{thebibliography}

\end{document}
